温度缩放(Temperature Scaling, TS)是深度分类器的默认事后重校准方法，然而其同分布(In-Distribution, ID)效益能否迁移到分布外(Out-of-Distribution, OOD)目标语料库，仍是跨机构心电图(Electrocardiogram, ECG)部署的核心开放问题。我们预注册并开展了一项多种子校准边界研究，覆盖 Chapman-Shaoxing、CPSC2018+2019 和 PTB-XL 之间的六个有向跨语料库迁移对，采用两种架构(InceptionTime 和 ResNet1D (1D-ResNet-34))、五个种子(42--46)、八种重校准方法，以及患者层级聚类配对自助法(operating CI 为 $B=10{,}000$ 的 BCa，在整个 60 实验网格上完全重跑；早期在 31/60 实验中 $B=10{,}000$ 的百分位 CI 和 21/60 实验中 $B=200$ 的结果作为历史溯源保留，两种 CI 方法在 60/60 实验中符号一致)。
预注册的主要终点是 OOD 校准效益 $\Delta\text{ECE}_{\text{OOD}}=\text{ECE}_{\mathrm{raw}}^{\mathrm{OOD}}
-\text{ECE}_{\mathrm{TS}}^{\mathrm{OOD}}$，针对已注册的经验方向先验进行单边检验。在 60 个种子实验(6 对 $\times$ 2 架构 $\times$ 5 种子，完全平衡)中，TS 在 51/60 (85.0\%) 个案例中产生正向 OOD 效益：InceptionTime 27/30 (90.0\%)，ResNet1D 24/30 (80.0\%)；九个反例以事后机理解释透明披露。当排除六个具有退化 CI 宽度($<3\times10^{-4}$)的实验后，标题计数变为 45/60 (75.0\%)，因此 85.0\% 的比率应与效应量和 CI 诊断联合解读。预注册的 ID 边界判据(95\% CI 包含零的单元 $\geq$80\%)\emph{未}满足(在 BCa operating CI 下 23/60 = 38.3\%；早期百分位 CI 给出 29/60 = 48.3\%)，触发预注册的``ID 非零边界''分支：OOD 效益约为 ID 效益的 $1.9\times$(均值 $+0.0159$ vs.\ $+0.0083$)，不能完全归因于 OOD 特有的失配。27 单元 Shapley 分解在合成数据($n=20{,}000$)上得到验证，方法选择边界将 TS 与最弱方法(EM 先验和矩阵缩放(Matrix scaling))分离，可预测性实验(LOO $R^2$ CI 上界 $0.104<0.5$)触发预注册的失败分支。判别感知校准门(C5)将校准有效性与临床可部署性解耦：7/14 对$\times$架构分层同时通过校准层和判别层。所有假设标记为\emph{探索性 + 鲁棒性验证}，而非验证性。综上，我们贡献 (i) 首次以预注册终点对 OOD 校准效益进行多种子量化；(ii) ID 边界条件分析并触发预注册分支；(iii) 在 27 个合成单元上的 Shapley 分解验证；(iv) 将 TS 与脆弱先验校正方法分离的方法选择边界；(v) 判别感知部署门，将 9 个反例重构为校准失败 vs.\ 判别失败。
